% CS5491: Reasoning Augmented FunSearch
% Author: 王宇浩 (72655080), 朱泳颐 (72657023), 李逸洋 (72655555)
% Responsibility: Shared project setup, analysis and report writing

% Compile with XeLaTeX; see scripts/build_milestone.py.
\documentclass[12pt,a4paper]{article}
\usepackage[margin=2.5cm]{geometry}
\usepackage{fontspec}
\usepackage{xeCJK}
\usepackage{amsmath}
\usepackage{booktabs,tabularx,array}
\usepackage{enumitem}
\usepackage{titlesec}
\usepackage{xcolor}
\usepackage[hidelinks]{hyperref}
\usepackage{fancyhdr}
\setmainfont{Times New Roman}
\setCJKmainfont{Songti SC}
\setmonofont{Menlo}[Scale=MatchLowercase]
\setCJKmonofont{Songti SC}
\setlength{\parindent}{0pt}
\setlength{\parskip}{7pt}
\linespread{1}
\raggedright
\titleformat{\section}{\normalfont\normalsize\bfseries}{}{0pt}{}
\titlespacing*{\section}{0pt}{12pt}{5pt}
\setlist[enumerate]{leftmargin=1.6em,topsep=3pt,itemsep=4pt,parsep=0pt}
\setlist[itemize]{leftmargin=1.4em,topsep=3pt,itemsep=4pt,parsep=0pt}
\renewcommand{\arraystretch}{1.35}
\newcolumntype{Y}{>{\raggedright\arraybackslash}X}
\pagestyle{fancy}
\fancyhf{}
\renewcommand{\headrulewidth}{0pt}
\fancyfoot[L]{\textcolor{gray}{CS5491 Artificial Intelligence}}
\fancyfoot[R]{\thepage}
\setlength{\footskip}{24pt}
\hypersetup{pdftitle={Reasoning Augmented FunSearch: Project Milestone},pdfauthor={Wang Yuhao; Zhu Yongyi; Li Yiyang}}

\begin{document}
\begin{center}
\textbf{Reasoning Augmented FunSearch for Online Bin Packing}\\[5pt]
Project Milestone\quad\textbar\quad 17 November 2026
\end{center}
\vspace{3pt}
\textbf{Group:} 王宇浩 (72655080), 朱泳颐 (72657023), 李逸洋 (72655555)\\[3pt]
\textbf{Contact:} \href{mailto:72655080@cityu-dg.edu.cn}{72655080@cityu-dg.edu.cn}

\section*{1. Problem and motivation}
In online bin packing, items arrive one at a time and must be placed immediately into bins with a fixed capacity. The aim is to use as few bins as possible. A decision that fills a bin well now may leave awkward gaps for later items, so simple rules do not always give the best packing.

Our project uses FunSearch-style program search to improve these rules. We ask: \textbf{Does giving the language model a short explanation of a packing strategy help it find better rules than giving it code alone?} A description might help the model develop an idea across several revisions, rather than only change the formula in the current program.

\section*{2. Method}
We compare two versions of the same search process:
\begin{itemize}
\item \textbf{Code only:} the model receives an existing scoring function and its performance, then proposes a revised function.
\item \textbf{Strategy plus code:} the model also receives a description of the current strategy, and returns a new description of up to 80 words together with the revised function.
\end{itemize}
Both versions use DeepSeek V4.1 Flash through the API, with the same model settings, starting programs and number of calls. First Fit and Best Fit provide the initial rules and reference results. First Fit chooses the first bin that can hold an item; Best Fit chooses the bin that leaves the smallest gap.

Each generated function scores the bins that can hold the current item. The evaluator chooses the highest score and opens a new bin only when none of the existing bins can fit the item. Programs are checked before evaluation, and those that run correctly are tested on the search instances. The search keeps a population of up to 20 programs and selects promising ones for the next revision.

\section*{3. Data and evaluation}
We prepared 100 search instances, 50 validation instances and 100 test instances. Each contains 500 items, with bin capacity 100. Item sizes come from a Weibull distribution with shape 3 and scale 45, rounded to integers and sampled again when outside 1--100. This is our own synthetic setup, informed by the released FunSearch examples.

We also prepared the 20 original instances in OR-Library's \texttt{binpack3} collection. They are divided into groups of 8, 4 and 8 before making three fixed item orders for each instance. This lets us examine whether rules found on the synthetic data also work on a different benchmark.

\newpage
\section*{4. Progress and preliminary results}
The evaluator, datasets, classical rules and both search versions are working. Runs can be saved and continued. The automated tests cover packing decisions, generated programs and the search process.

A short pilot made eight API calls per method. All 16 generated programs ran successfully, confirming that the model connection and evaluation loop worked together.

\begin{center}
\begin{tabularx}{\linewidth}{@{}Yrrr@{}}
\toprule
\textbf{Pilot method} & \textbf{Calls} & \textbf{Valid} & \textbf{Best ratio}\\
\midrule
Code only & 8 & 8 & 1.000\\
Strategy plus code & 8 & 8 & 1.000\\
\bottomrule
\end{tabularx}
\end{center}
\vspace{-2pt}
For each instance, we divide each program's bin count by Best Fit's count, then average across instances. A ratio of 1 matches Best Fit; lower is better. Both pilots ended at 1.000. The pilot checked the setup; the full comparison examines whether the strategy helps across repeated runs.

On the 100 search instances, Best Fit used 21,155 bins and First Fit used 21,255. These results provide a reference for assessing generated rules.

\section*{5. Workload ownership}
\begin{tabularx}{\linewidth}{@{}p{4.1cm}Y@{}}
\toprule
\textbf{Group member} & \textbf{Area of responsibility}\\
\midrule
王宇浩 (72655080) & \textbf{LLM interface and strategy augmentation.} Owns the DeepSeek API connection, prompts, response handling and the strategy-plus-code version.\\[4pt]
朱泳颐 (72657023) & \textbf{Search framework and logging.} Owns program selection, population updates, the search loop and saved run records.\\[4pt]
李逸洋 (72655555) & \textbf{Benchmark preparation and evaluation scripts.} Owns dataset preparation, the packing evaluator and experiment scripts.\\
\bottomrule
\end{tabularx}

All three members share the analysis of results and report writing.

\section*{6. Full comparison and supporting code}
The full comparison uses five independent runs per method and 200 calls per run. Validation data chooses one program from each run; these are then compared on test data. We examine bin use, progress during search, the number of programs that work and API usage.

The code, evaluator and pilot records are available at:\\[3pt]
\href{https://github.com/alanwangcot/cs5491-funsearch/tree/main}{\texttt{github.com/alanwangcot/cs5491-funsearch}}\\[3pt]
Milestone code is on the \texttt{main} branch. The README gives setup and evaluation commands.
\end{document}
